% !TeX root = ../main.tex
\section{Working without the dry signal}
\label{sec:nodi}

Everything in Section~\ref{sec:modelling} needs the dry signal of the same
performance. A song does not have it. There are three ways around this.

\subsection{Unpaired and blind modelling}

\begin{keyidea}[: learn the device from its output's statistics]
If we cannot compare the model's output with the target sample by sample, we
can still ask that the model's output \emph{looks like} processed guitar in
general. A \term{discriminator} network trained to tell real processed guitar
from the model's output provides that signal (adversarial training, as in GANs).
\end{keyidea}

Wright, Välimäki and Juvela trained a real-time WaveNet amp model only
adversarially~\cite{wright2023adversarial}. With input and target from the same
guitar, the best unpaired model scored 83--92 MUSHRA across three plugin tones;
with a different guitar for the input, the realistic case, 77--84, comparable to
an oracle that equalised the true plugin's output (72--81) \reported. Training
took 4--10 GPU hours per tone. Chen et al.\ improved stability with multi-scale
and multi-period discriminators and extra clean data, and found loudness
normalisation essential~\cite{chen2024unsupervised}.

Moliner et al.\ estimate the clean signals and the effect jointly by
\term{diffusion posterior sampling} with a prior trained on clean
guitar~\cite{moliner2025unsupervised}. Diffusion stays stable with only 18\,s of
processed audio, where adversarial training degrades; adversarial training is
better on heavy distortion; a grey-box operator matched black-box ones under
diffusion but collapsed under adversarial training. Švento et al.\ estimate a
memoryless nonlinearity (a spline) together with the clean
signal~\cite{svento2025estimation}, giving both a dry estimate and a waveshaper
curve; LOUDAR moves the operator into an autoencoder's latent
space~\cite{svento2026loudar}.

All of these were evaluated on isolated, plugin-rendered guitar, mostly with
distortion alone. None handles cabinet, modulation, delay and reverb together,
and none was tested on a separated stem.

\subsection{Effect removal}

A recovered dry signal turns the wet-only problem back into a paired one. Imort
et al.\ removed distortion with separation networks~\cite{imort2022distortion}.
RemFX detects which of five effects are present and chains effect-specific
removers; SI-SDR falls from 16.7\dB{} with one effect to 2.6\dB{} with
five~\cite{rice2023general}. Distortion Recovery restores a mel spectrogram and
re-synthesises it with a vocoder~\cite{lee2024distortion}. EG-VAE splits a
processed recording into frame-level content and a global tone embedding;
swapping the embedding transfers the tone, a masking mode removes
it~\cite{yeh2026egvae}. On plugin tones it reaches a tone-similarity opinion
score of 4.30 on seen and 3.70 on unseen tones (oracle 4.48; baselines at most
2.74), and a ``dryness'' score of 3.42--3.50 for removal against 4.39--4.58 for
the true DI \reported. A baseline that fitted a fixed differentiable chain of
generic blocks to EG-VAE's dry estimate scored only 1.44--2.74 in tone
similarity: evidence that generic EQ, compressor and distortion blocks do not
cover guitar tone diversity. No code or weights were found for EG-VAE.

Removal and remodelling have been chained: Hinrichs et al.\ remove a single
effect with a separation network, then train a model on the recovered pair;
listeners rated the copies at medians of 92--100 for slapback, phaser and soft
clipping~\cite{hinrichs2024blind}. SDR does not track perceived dryness for
distortion, as also found by Imort et al.\ and the EG-VAE authors.

\subsection{Dry signals from other sources}

Three other sources of a dry signal exist, each with its own mismatch:
\begin{center}
\small
\begin{tabular}{@{}>{\raggedright}p{4.2cm}ccc@{}}
\toprule
Source & Same notes and timing & Same guitar timbre & Artefacts \\
\midrule
A different performance (corpus DI, the player's DI) & no & roughly & none \\
Resynthesis of the transcribed notes & yes & no & transcription errors \\
An effect-removal model & yes & yes & removal errors \\
\bottomrule
\end{tabular}
\end{center}
A different performance supports only unaligned losses, as in unpaired
modelling~\cite{wright2023adversarial}. Resynthesis gives the same notes at the
same times with the wrong timbre; GuitarFlow turns sample-based renders of
tablature into realistic DI with flow matching~\cite{loth2025guitarflow}, but for
data generation, not tone estimation. \textbf{We found no published work that
resynthesises a DI from transcribed notes in order to fit a tone.} That is the
route the app takes first (Section~\ref{sec:impl-dry}), because the app already
transcribes the notes.
